\documentclass[a4paper]{article}
\usepackage[PUTF,T1]{fontenc}
\usepackage[utf8]{inputenc}
\usepackage{tabularx}
\usepackage{longtable}
\usepackage[table]{xcolor}
\usepackage{graphicx}
\usepackage{datetime}
\usepackage{libertine}
\usepackage{ngerman}
\usepackage{fancyhdr}
\usepackage{lastpage}
\usepackage{xpatch}
\usepackage{seqsplit}

\usepackage{avant} %use helvetica?


\usepackage{geometry}
\geometry{
  a4paper,
  total={170mm,257mm},
  left=20mm,
  top=35mm,
  bottom=35mm,
  heightrounded,
 }
 
 %Thanks, Ulrike: https://github.com/u-fischer/putfenc
 \makeatletter

% Unterdrücken der aux-Ausgabe:
\def\qr@writebinarymatrixtoauxfile#1{}%
\makeatother

\DeclareFontFamily{PUTF}{pdf}{}%
\DeclareFontShape{PUTF}{pdf}{m}{n}{ <-> ecrm1000 }{}%
\DeclareFontSubstitution{PUTF}{pdf}{m}{n}%


\definecolor{papercolor}{rgb}{1,1,1}
\definecolor{bgcolor}{rgb}{1,1,1}
\definecolor{textcolor}{rgb}{0,0,0}

\pagecolor{bgcolor}


\usepackage[xetex,%
            colorlinks=true,linkcolor=blue,citecolor=blue,%
            anchorcolor=red,urlcolor=blue,bookmarks=true,%
            bookmarksopen=true,bookmarksopenlevel=0,plainpages=false,%
            bookmarksnumbered=true,hyperindex=false,pdfstartview=,%
            pdfauthor={StuStaPay},%
            pdftitle={StuStaPay Umsatzbericht},%
            %pdfsubject={StuStaPay Umsatzbericht}%
]{hyperref}


% for deterministic build
\special{pdf:trailerid [
    <00112233445566778899aabbccddeeff>
    <00112233445566778899aabbccddeeff>
]}

\begin{document}
    \color{textcolor}

    \newgeometry{
      a4paper,
      total={170mm,257mm},
      left=20mm,
      top=0mm,
      bottom=35mm,
      heightrounded,
    }

    \pagestyle{fancy}
    \fancypagestyle{plain}{ % pages containing chapter start
        \fancyhf{}
        \fancyhead{} % clear all header fields
        \renewcommand{\headrulewidth}{0pt}
    }
    \fancyfoot[C]{\includegraphics[width=3cm]{logo} \\\textcolor{textcolor}Seite \thepage/\pageref*{LastPage}}


    \noindent
    \makebox[17cm]{
    \vspace{-35mm}
    \includegraphics[width=21cm,trim={15.5cm 0 10.15cm 0},clip]{stustaculum24_wristband.pdf}
    }
    \vspace{1em}

    \begin{center}
        \LARGE\textsf{StuStaPay} \\
        \Huge{Umsatzbericht} \\
        \LARGE\textsf{StuStaCulum 2024}
    \end{center}
    \vspace{2em}

    % Addresse und Steuer ID
    \begin{flushright}
        \VAR[config.issuer|latex] \\
        \VAR[config.address.strip()|latex] \\
        USt-IdNr. \VAR[config.ust_id|latex] \\
    \end{flushright}

    \noindent
    \VAR[node.name] \\
    \VAR[node.description] \\

    \vspace{2em}

    \begin{center}
        \LARGE\textbf{Übersicht} \\
    \end{center}
    \vspace{2em}

    {
    \large
    \begin{tabular*}{10cm}{ p{6cm} p{6cm} }
        Beginn                             & \VAR[from_time|datetime] \\
        Letzter berücksichtigter Zeitpunkt & \VAR[to_time|datetime] \\
        Gesamtumsatz                       & \VAR[total_revenue|money] \\
        Gebühren                           & \VAR[fees|money] (\VAR[fees_percent|percent]) \\
        Auszahlung nach Gebühren           & \VAR[revenue_minus_fees|money] \\
    \end{tabular*}
    }

    \vspace{2em}

    \begin{center}
        \LARGE\textbf{Nach Tagen} \\
    \end{center}
    \vspace{2em}
    
    % Auflistung nach Tagen
    \newcolumntype{R}{>{\raggedleft\arraybackslash}X}
    \rowcolors{2}{}{bgcolor!80}
    \begin{tabularx}{\textwidth}{ R R R R R }
        \textbf{Datum} & \textbf{Anzahl Transaktionen} & \textbf{Gesamtumsatz} & \textbf{Gebühren} & \textbf{Auszahlung} \\
        \hline
        \BLOCK[ for daily in daily_revenue_stats ]
        \VAR[daily.day] & \VAR[daily.num_transactions] & \VAR[daily.revenue|money] & \VAR[daily.fees|money] & \VAR[daily.revenue_minus_fees|money] \\
        \BLOCK[endfor]
    \end{tabularx}

    \vfill
    \restoregeometry
    \pagebreak

    % hier ist null abstand oben.
    % vermutlich wegen dem fancy header, hatte keine zeit mehr das zu fixen
    % null abstand oben haben wir auch bei natürlichen page brakes
    % muss korrigiert werden

    \begin{center}
        \LARGE\textbf{Einzelauflistung} \\
    \end{center}
    \vspace{1em}

    \rowcolors{0}{}{}

    \begin{center}
    \begin{longtable}{lllll}
    \hline \multicolumn{1}{c}{\textbf{Datum \& Uhrzeit}} & \multicolumn{1}{c}{\textbf{Transaktionsnummer}} & \multicolumn{1}{c}{\textbf{Betrag}} & \multicolumn{1}{c}{\textbf{Gebühr}} & \multicolumn{1}{c}{\textbf{Auszahlung}} \\ \hline 
    \endfirsthead
    
    \hline
    \multicolumn{5}{l}
    {{... Fortsetzung von der vorherigen Seite}} \\
    \hline \multicolumn{1}{c}{\textbf{Datum \& Uhrzeit}} & \multicolumn{1}{c}{\textbf{Transaktionsnummer}} & \multicolumn{1}{c}{\textbf{Betrag}} & \multicolumn{1}{c}{\textbf{Gebühr}} & \multicolumn{1}{c}{\textbf{Auszahlung}} \\ \hline
    \endhead
    
    \hline \multicolumn{5}{r}{{Fortsetzung auf der nächsten Seite ...}} \\ \hline
    \endfoot
    
    \hline \hline
    \endlastfoot
        \BLOCK[ for order in orders ]
        \VAR[order.booked_at|datetime] & \VAR["{:010}".format(order.id)] & \VAR[order.total_price|money] & \VAR[order.fees|money] & \VAR[order.total_price_minus_fees|money] \\
        \BLOCK[endfor]
    \end{longtable}
    \end{center}

    \begin{tabularx}{\textwidth}{ RRRRR }
    \end{tabularx}

\end{document}
